\section{Polarización dodecafásica, dualidad reticular y reciprocidad theta}
\label{sec:leech-dualidad}

La incidencia excepcional y la polarización del registro dodecafásico
proceden de dos realizaciones del mismo estado enriquecido. La primera
conserva las relaciones de frontera, el código de pegado y el origen
marcado que determinan el retículo; la segunda hace actuar una
representación sobre las doce coordenadas enteras ya obtenidas por los
lectores del Topological Phase Kernel. Su composición permite estudiar
una deformación métrica cuyo parámetro modifica las longitudes relativas
de doce planos y conserva simultáneamente el volumen total y una
simetría ternaria explícita.

El punto de partida causal sigue siendo la construcción
APP--TRIT--TPK del estado enriquecido y del continuo conjunto.
El registro \(K\) se recibe después del registro dodecafásico y de sus
operaciones de incidencia y orientación. La representación empleada
actúa sobre esa salida y conserva sus doce valores como datos ya
determinados. La realización reticular se recibe con su forma bilineal,
sus clases de pegado y su vector marcado. Estas estructuras, desarrolladas
respectivamente en \cite{hmtX,HMT-IV}, se mantienen al componerlas.
La colección que se construye a continuación es una realización geométrica
de esos datos con una coordenatización común de doce posiciones.

\subsection{Polarización del registro y marcación de los doce planos}
\label{subsec:leech-polarizacion}

El registro dodecafásico y su suma son
\begin{equation}
 \begin{split}
 K={}&(234,543,140,729,659,824,\\
     &\hspace{13mm}621,58,914,794,146,601),\\
 \sum_{j=1}^{12}K_j={}&6263.
 \end{split}
 \label{eq:leech-K}
\end{equation}
Sea \(A_5\) el grupo de permutaciones pares de cinco elementos y sea
\(C_5=\langle(1\,2\,3\,4\,5)\rangle\). Sus doce clases izquierdas
determinan una representación por permutaciones
\(\varrho:A_5\to O(\mathbb R^{12})\). La coordenatización marcada se fija mediante
los representantes siguientes, escritos como listas de imágenes:
\[
\begin{array}{c|c@{\qquad}c|c}
j&r_j&j&r_j\\ \hline
1&(4,3,5,2,1)&7&(5,4,3,2,1)\\
2&(4,2,1,3,5)&8&(1,3,4,2,5)\\
3&(1,2,4,5,3)&9&(4,2,3,5,1)\\
4&(2,5,3,4,1)&10&(3,2,5,4,1)\\
5&(4,3,1,5,2)&11&(4,5,2,3,1)\\
6&(5,1,2,3,4)&12&(5,3,2,4,1).
\end{array}
\]
Concretamente, \(\varrho(a)e_j=e_\ell\) si
\(ar_jC_5=r_\ell C_5\). El carácter tridimensional \(\chi_3\) toma
los valores
\[
\begin{array}{c|ccccc}
\text{clase}&1&2^2&3&5_a&5_b\\ \hline
\chi_3&3&-1&0&(1+\sqrt5)/2&(1-\sqrt5)/2.
\end{array}
\]
La clase \(5_a\) contiene el generador de \(C_5\) indicado y \(5_b\)
contiene su cuadrado. Esta elección distingue los dos caracteres
tridimensionales conjugados. Si \(I_{12}\) es la identidad y
\(\mathbf1\) es el vector de doce unos, se definen
\begin{equation}
 P_3=\frac3{60}\sum_{a\in A_5}\chi_3(a^{-1})\varrho(a),
 \qquad
 P_{11}=I_{12}-\frac1{12}\mathbf1\mathbf1^{\mathsf T}.
 \label{eq:leech-projectors}
\end{equation}
El campo \(\mathbb Q(\sqrt5)\) aparece aquí en la realización de la
representación que actúa sobre el registro generado.

\begin{lemma}[Polarización de suma nula]
\label{lem:leech-u}
Los operadores de \eqref{eq:leech-projectors} satisfacen
\[
 P_3=P_3^{\mathsf T}=P_3^2,\qquad
 \operatorname{rank}P_3=3,\qquad P_3\mathbf1=0,\qquad
 P_3P_{11}=P_{11}P_3=P_3.
\]
Por consiguiente,
\begin{equation}
 u=P_3P_{11}K=P_3K,\qquad
 \sum_{j=1}^{12}u_j=0,\qquad
 \|u\|^2=\frac{6638585+2275584\sqrt5}{20}>0.
 \label{eq:leech-u-norm}
\end{equation}
\end{lemma}
\begin{proof}
La suma del carácter es el proyector isótipo de la representación.
La realidad del carácter y la sustitución \(a\mapsto a^{-1}\)
dan su simetría. La multiplicidad del carácter tridimensional en la
representación de \(A_5/C_5\) es la dimensión de sus invariantes por
\(C_5\), que vale
\[
 \frac15\left(
 3+2\frac{1+\sqrt5}{2}+2\frac{1-\sqrt5}{2}
 \right)=1.
\]
El proyector tiene por ello rango tres. Su ortogonalidad al carácter
trivial implica \(P_3\mathbf1=0\) y las identidades con \(P_{11}\).
La aplicación de la suma finita al registro
\eqref{eq:leech-K} produce las doce coordenadas
\begin{equation}
\begin{aligned}
u={}&(-495/4-233\sqrt5/10,\quad403/4+169\sqrt5/10,\\
&-403/4-169\sqrt5/10,\quad495/4+233\sqrt5/10,\\
&601/4+1031\sqrt5/10,\quad203/4+211\sqrt5/5,\\
&-203/4-211\sqrt5/5,\quad-601/4-1031\sqrt5/10,\\
&313/4+569\sqrt5/10,\quad162+219\sqrt5/20,\\
&-162-219\sqrt5/20,\quad-313/4-569\sqrt5/10).
\end{aligned}
\label{eq:leech-u-coordinates}
\end{equation}
Sumarlas y elevarlas al cuadrado en \(\mathbb Q(\sqrt5)\) da
\eqref{eq:leech-u-norm}. Equivalentemente,
\[
 \|u\|^2=K^{\mathsf T}P_3K
 =\frac1{20}\sum_{a\in A_5}
          \chi_3(a^{-1})K^{\mathsf T}\varrho(a)K.
\]
Los coeficientes positivos de su evaluación prueban que \(u\ne0\).
\end{proof}

La suma nula tiene una función geométrica concreta: las expansiones de
unos planos podrán compensar las contracciones de los restantes en el
determinante total. La no nulidad de \(u\), en cambio, garantiza que esa
compensación contiene una deformación efectiva.

El espacio euclídeo ambiente de la realización reticular es
\begin{equation}
 E=\bigoplus_{j=1}^{12}E_j,\qquad
 B=\operatorname{diag}(B_2,\ldots,B_2),\qquad
 B_2=\begin{pmatrix}2&-1\\-1&2\end{pmatrix}.
 \label{eq:leech-ambient}
\end{equation}
Cada \(E_j\) es el plano real generado por el retículo de raíces
\(A_2\), expresado en su base de raíces simples. La descomposición
ortogonal pertenece al espacio ambiente. El retículo de Leech se
obtiene mediante pegado y paso a un vecino, operaciones que relacionan
los doce factores.

Sea \(C_W\subset\mathbb F_3^{12}\) el código ternario de dimensión seis
de la construcción Paley--Witt. Identificando el discriminante de cada
factor \(A_2\) con \(\mathbb F_3\), el pegado es
\[
 N=N(C_W)=
 \bigcup_{c\in C_W}\bigl(A_2^{12}+\phi(c)\bigr),
\]
donde \(\phi(c)\) elige representantes de las clases discriminantes.
La autoortogonalidad y autodualidad del código transmiten paridad e
integralidad a \(N\); su determinante es
\[
 \det N=\frac{3^{12}}{(3^6)^2}=1.
\]
La construcción excepcional fija el origen marcado \(j=1\) y el
vector
\begin{equation}
 v_\alpha=4\rho_1+\rho_2+\cdots+\rho_{12},
 \qquad \rho_j=(1,1)_j,\qquad
 \|v_\alpha\|_B^2=54.
 \label{eq:leech-marked-vector}
\end{equation}
Con él se forman
\begin{equation}
\begin{split}
 N_{\alpha,0}
  &=\{x\in N:\langle x,v_\alpha\rangle_B\equiv0\pmod3\},\\
 \Lambda&=N_{\alpha,0}+\mathbb Z\,v_\alpha/3.
\end{split}
\label{eq:leech-neighbour}
\end{equation}
La realización excepcional establece que \(\Lambda\) es un retículo
positivo, par, unimodular, sin raíces y de rango \(24\), de modo que es
isométrico al retículo de Leech \cite{HMT-IV}. En particular,
\(\Lambda^*=\Lambda\), con dual respecto de \(B\).
Estas propiedades y la pertenencia \(v_\alpha/3\in\Lambda\) son las
hipótesis reticulares precisas empleadas en los teoremas siguientes.

\subsection{Isometría ternaria y estabilidad del vecino}
\label{subsec:leech-c3}

Definimos las matrices bidimensionales
\[
 C=\begin{pmatrix}0&-1\\1&-1\end{pmatrix},
 \qquad
 R=\begin{pmatrix}0&1\\1&0\end{pmatrix},
 \qquad
 \varpi=\frac13\begin{pmatrix}2\\1\end{pmatrix}.
\]
La multiplicación directa verifica
\[
 C^3=I_2,\quad C^2+C+I_2=0,\quad
 C^{\mathsf T}B_2C=B_2,\quad
 RCR=C^{-1},\quad R^{\mathsf T}B_2R=B_2.
\]
Además, \(C\varpi-\varpi=(-1,0)^{\mathsf T}\), por lo que \(C\)
actúa trivialmente en \(A_2^*/A_2\). La reflexión \(R\) cambia el signo
de la clase discriminante. El código \(C_W\) contiene la palabra de
soporte completo
\[
 c_*=(1,1,1,1,1,1,2,1,1,1,1,1).
\]
En la presentación \(C_W=\{(q,qA_W):q\in\mathbb F_3^6\}\), dicha
palabra procede de \(q_*=(1,1,1,1,1,1)\) y
\(q_*A_W=(2,1,1,1,1,1)\).

Para cada plano se fija el marco
\begin{equation}
 Q_j=\begin{cases}
 R,&(c_*)_j=1,\\
 I_2,&(c_*)_j=2,
 \end{cases}
 \qquad
 g_c=\bigoplus_{j=1}^{12}Q_jCQ_j^{-1}.
 \label{eq:leech-c3}
\end{equation}
Los símbolos \(Q_j\) designan marcos de planos; el símbolo \(P_3\)
queda reservado para el proyector sobre el sector tridimensional.

\begin{proposition}[Simetría de orden tres del retículo marcado]
\label{prop:leech-c3}
La transformación \(g_c\) es una isometría de orden tres de \(E\),
tiene como único vector fijo el cero y satisface
\[
 g_cN=N,\qquad g_cN_{\alpha,0}=N_{\alpha,0},
 \qquad g_c\Lambda=\Lambda.
\]
\end{proposition}
\begin{proof}
Cada bloque es \(C\) o \(C^{-1}\), cuyo polinomio mínimo es
\(X^2+X+1\). Esto prueba el orden y la afirmación sobre vectores fijos,
así como la conservación de \(B\). Cada bloque actúa trivialmente en
el discriminante y conserva, por tanto, todas las clases de pegado de
\(N\).

Escribamos \(v=v_\alpha=\sum_j a_j\rho_j\), con \(a_1=4\) y
\(a_j=1\) para \(j>1\). Todos los \(a_j\) son congruentes con uno
módulo tres. Las identidades
\[
 (C-I_2)\rho=(-2,-1)^{\mathsf T},
 \qquad (C^{-1}-I_2)\rho=(-1,-2)^{\mathsf T}
\]
muestran que
\[
 y=\frac{g_cv-v}{3},\qquad z=\frac{g_c^{-1}v-v}{3}
\]
tienen clases discriminantes \(c_*\) y \(-c_*\), respectivamente.
Ambas pertenecen a \(C_W\), luego \(y,z\in N\). En cada factor,
\[
 \left\langle
 \frac{(C^{\pm1}-I_2)a_j\rho}{3},a_j\rho
 \right\rangle_{B_2}=-a_j^2.
\]
Por suma, \(\langle y,v\rangle_B=\langle z,v\rangle_B=-27\);
por ello \(y,z\in N_{\alpha,0}\).
Si \(x\in N_{\alpha,0}\), entonces
\[
 \langle g_cx,v\rangle_B
  =\langle x,g_c^{-1}v\rangle_B
  =\langle x,v\rangle_B+3\langle x,z\rangle_B
  \equiv0\pmod3.
\]
La misma operación con \(g_c^{-1}\) demuestra la igualdad
\(g_cN_{\alpha,0}=N_{\alpha,0}\). Finalmente,
\(g_c(v/3)=v/3+y\) conserva la clase que genera
\(\Lambda=N_{\alpha,0}+\mathbb Z(v/3)\).
\end{proof}

La simetría ternaria actúa, por consiguiente, sobre el retículo
pegado completo. Su conservación utiliza la palabra del código y
la congruencia del vector marcado, además de la rotación de cada
plano. Esta precisión permite transportar la simetría al flujo sin
perder la información aritmética que hizo posible su construcción.

\subsection{Flujo anisótropo y conservación del covolumen}
\label{subsec:leech-flow}

\begin{definition}[Deformación determinada por la polarización]
Para \(s\in\mathbb R\), parámetro adimensional, se define
\begin{equation}
 H_u=\bigoplus_{j=1}^{12}u_jI_2,\qquad
 D_s=\exp(sH_u)=\bigoplus_{j=1}^{12}e^{su_j}I_2,\qquad
 \Lambda_s=D_s\Lambda.
 \label{eq:leech-flow}
\end{equation}
La coordenada \(u_j\) se asigna al plano \(E_j\) de la misma coordenatización
marcada.
\end{definition}

El factor positivo \(e^{su_j}\) conserva la orientación de cada
plano. Las doce etiquetas transmiten la correspondencia entre la
polarización y la realización reticular. Esta correspondencia
forma parte de la estructura de partida; una permutación de las
etiquetas se trata mediante su transporte simultáneo en ambos objetos.

\begin{theorem}[Covolumen y simetría del flujo]
\label{thm:leech-flow}
Para todos \(s,t\in\mathbb R\) se cumple
\[
\begin{gathered}
 D_{s+t}=D_sD_t,\qquad D_0=I,\qquad D_s^{-1}=D_{-s},\\
 D_s^\dagger=D_s,\qquad D_sg_c=g_cD_s,\qquad \det D_s=1,
\end{gathered}
\]
donde \(A^\dagger=B^{-1}A^{\mathsf T}B\).
Cada \(\Lambda_s\) es una red discreta de rango \(24\) y covolumen
uno, preservada por la isometría \(g_c\) de orden tres.
\end{theorem}
\begin{proof}
Las identidades de grupo se verifican en cada bloque exponencial.
Los bloques escalares son autoadjuntos para \(B_2\), y conmutan
con los bloques de \(g_c\). La suma nula de
\eqref{eq:leech-u-norm} da
\[
 \det D_s=\prod_{j=1}^{12}e^{2su_j}
 =\exp\left(2s\sum_{j=1}^{12}u_j\right)=1.
\]
La imagen de una red por un operador invertible es una red del
mismo rango; su covolumen se multiplica por el valor absoluto del
determinante. El covolumen de \(\Lambda\) es uno, de modo que el
de \(\Lambda_s\) también lo es. La conmutación y
la proposición~\ref{prop:leech-c3} dan
\[
 g_c\Lambda_s=g_cD_s\Lambda=D_sg_c\Lambda=D_s\Lambda.
\]
El orden y el espacio de puntos fijos de \(g_c\) permanecen
inalterados.
\end{proof}

\begin{remark}[Parámetro y simetría conservada]
La dirección \(u\) queda determinada por el registro; \(s\) recorre
una colección de realizaciones. Reemplazar \(u\) por \(u/\|u\|\)
reparametriza la misma colección. Una interpretación física que
seleccione un valor particular de \(s\) añade su regla de realización.
El teorema conserva la simetría explícita \(g_c\); el transporte de
otros automorfismos requiere sus correspondientes relaciones con
\(H_u\).
\end{remark}

\subsection{Dual euclídeo y deformación efectiva}
\label{subsec:leech-dual}

\begin{definition}[Dual euclídeo]
Para una red \(L\subset E\), su dual respecto de \(B\) es
\[
 L^*=\{y\in E:\langle y,x\rangle_B\in\mathbb Z
                         \text{ para todo }x\in L\}.
\]
\end{definition}

\begin{theorem}[Reciprocidad del dual]
\label{thm:leech-dual}
La colección \eqref{eq:leech-flow} satisface
\begin{equation}
 \Lambda_s^*=\Lambda_{-s}\qquad(s\in\mathbb R).
 \label{eq:leech-dual}
\end{equation}
Su autodualidad en un parámetro \(s\) equivale a
\(D_{2s}\Lambda=\Lambda\).
\end{theorem}
\begin{proof}
Para un operador invertible \(A\),
\[
 (AL)^*=A^{-\dagger}L^*.
\]
En efecto, \(\langle y,Ax\rangle_B=\langle A^\dagger y,x\rangle_B\)
es entero para todo \(x\in L\) precisamente cuando
\(A^\dagger y\in L^*\). Tomando \(A=D_s\), la autoadjunción,
la identidad \(D_s^{-1}=D_{-s}\) y \(\Lambda^*=\Lambda\) prueban
\eqref{eq:leech-dual}. La igualdad
\(D_s\Lambda=D_{-s}\Lambda\) equivale, multiplicando por \(D_s\),
a \(D_{2s}\Lambda=\Lambda\).
\end{proof}

El covolumen controla el volumen de un dominio fundamental; la
autodualidad controla todos los productos enteros con la red.
La primera propiedad es continua a lo largo de esta colección. La
segunda impone la condición aritmética exacta del teorema anterior.
Se empleará por ello la expresión \emph{covolumen uno} para el
flujo euclídeo y se especificará la integralidad cuando intervenga
la unimodularidad aritmética.

\begin{proposition}[Testigo de pérdida local de integralidad]
\label{prop:leech-nonintegral}
Existe \(\delta>0\) tal que, si \(0<|s|<\delta\), la red
\(\Lambda_s\) es no integral y no es isométrica al retículo de Leech.
\end{proposition}
\begin{proof}
El vector \(v_\alpha/3\) pertenece a \(\Lambda\). Su norma
transportada es
\begin{equation}
 f(s)=\left\|D_s\frac{v_\alpha}{3}\right\|_B^2
 =\frac29\left(16e^{2su_1}+
                       \sum_{j=2}^{12}e^{2su_j}\right).
 \label{eq:leech-witness}
\end{equation}
Por \(\sum_j u_j=0\) y por la primera coordenada de
\eqref{eq:leech-u-coordinates},
\begin{equation}
\begin{split}
 f(0)&=6,\\
 f'(0)&=\frac49\left(16u_1+\sum_{j=2}^{12}u_j\right)
       =\frac{20}{3}u_1
       =-825-\frac{466}{3}\sqrt5\ne0.
\end{split}
\label{eq:leech-witness-derivative}
\end{equation}
La continuidad de \(f'\) permite elegir un intervalo
\((-\delta,\delta)\) en el que \(f\) sea estrictamente monótona.
Reduciendo \(\delta\), se obtiene además
\(11/2<f(s)<13/2\). Para \(0<|s|<\delta\), resulta \(f(s)\ne6\);
en ese intervalo, seis es el único entero. Por tanto, \(\Lambda_s\)
contiene un vector de norma no entera. Una red integral tiene
norma entera en todos sus vectores. Puesto que Leech es integral
y la norma se conserva por isometrías, queda demostrada también
la segunda afirmación.
\end{proof}

La deformación preserva volumen y simetría ternaria al tiempo que
modifica efectivamente la estructura métrica aritmética. Los pares
opuestos de coordenadas de \(u\) definen una permutación del espacio
ambiente que intercambia dilataciones y contracciones. Su promoción
a un automorfismo del retículo pegado exigiría conservar el código
y el vecino. La reciprocidad del dual ya demostrada utiliza
únicamente la autoadjunción del flujo y la autodualidad inicial.

\subsection{Convergencia y reciprocidad de las funciones theta}
\label{subsec:leech-theta}

Para \(s\in\mathbb R\) y \(t>0\), definimos
\begin{equation}
 \Theta_s(t)=\sum_{\lambda\in\Lambda}
                \exp\bigl(-\pi t\|D_s\lambda\|_B^2\bigr).
 \label{eq:leech-theta-real}
\end{equation}
Es la suma gaussiana de todas las longitudes de la red deformada,
con sus multiplicidades. Si \(M=\max_j|u_j|\), se tiene
\begin{equation}
 e^{-2|s|M}\|x\|_B^2
 \leq\|D_sx\|_B^2
 \leq e^{2|s|M}\|x\|_B^2.
 \label{eq:leech-quadratic-bounds}
\end{equation}
Sobre un compacto de \(\mathbb R\times(0,\infty)\), la cota
inferior proporciona una gaussiana dominante sobre la red fija.
La serie converge absoluta y uniformemente. Las derivadas de
orden finito añaden factores polinomiales en las coordenadas de
\(\lambda\); dichos factores quedan dominados por otra gaussiana.
Las derivaciones término a término son, por consiguiente,
legítimas en esos compactos.

\begin{theorem}[Reciprocidad gaussiana]
\label{thm:leech-theta}
Con el volumen euclídeo asociado a \(B\),
\begin{equation}
 \Theta_s(t)=t^{-12}\Theta_{-s}(t^{-1}),
 \qquad s\in\mathbb R,\quad t>0.
 \label{eq:leech-theta-reciprocity}
\end{equation}
En particular, \(\Theta_s(1)=\Theta_{-s}(1)\).
\end{theorem}
\begin{proof}
Fijamos la transformada de Fourier
\[
 \widehat f(y)=\int_E
      f(x)e^{-2\pi i\langle x,y\rangle_B}\,dx.
\]
Para una red \(L\subset E\) y
\(f_t(x)=e^{-\pi t\|x\|_B^2}\), periodizamos
\[
 F_t(x)=\sum_{\ell\in L}f_t(x+\ell).
\]
La convergencia normal de la serie y sus derivadas da una función
suave sobre \(E/L\). Si \(\xi\in L^*\), su coeficiente de Fourier es
\[
\begin{split}
 c_\xi
 &=\frac1{\operatorname{covol}(L)}
   \int_{E/L}F_t(x)e^{-2\pi i\langle x,\xi\rangle_B}\,dx\\
 &=\frac{\widehat f_t(\xi)}{\operatorname{covol}(L)}.
\end{split}
\]
La segunda igualdad descompone \(E\) en traslaciones de un dominio
fundamental; el carácter tiene valor uno sobre \(L\). El producto
de las integrales gaussianas unidimensionales, en una base
ortonormal para \(B\), da
\[
 \widehat f_t(\xi)=
 t^{-12}\exp\bigl(-\pi\|\xi\|_B^2/t\bigr).
\]
La serie de Fourier es absolutamente convergente. Evaluándola en
\(x=0\), obtenemos la suma de Poisson
\[
 \sum_{\ell\in L}e^{-\pi t\|\ell\|_B^2}
 =\frac{t^{-12}}{\operatorname{covol}(L)}
      \sum_{\xi\in L^*}e^{-\pi\|\xi\|_B^2/t}.
\]
Ahora \(L=\Lambda_s\) tiene covolumen uno y
\(L^*=\Lambda_{-s}\) por el teorema~\ref{thm:leech-dual}.
La sustitución prueba \eqref{eq:leech-theta-reciprocity}.
El caso \(t=1\) es inmediato.
\end{proof}

La periodización anterior es la demostración gaussiana de Poisson;
puede cotejarse con \cite[teorema 2, pp.~10--11]{Elkies}.
La convención con signo positivo en Fourier conduce a la misma
transformada de la gaussiana, por su paridad. La igualdad de theta
afecta a las series completas y conserva el dominio analítico de
convergencia.

\begin{proposition}[Extensión holomorfa]
\label{prop:leech-theta-complex}
Para \(\tau\) en el semiplano superior
\(\mathbb H=\{\tau\in\mathbb C:\operatorname{Im}\tau>0\}\), la serie
\begin{equation}
 \theta_s(\tau)=\sum_{\lambda\in\Lambda}
             e^{\pi i\tau\|D_s\lambda\|_B^2}
 \label{eq:leech-theta-complex}
\end{equation}
converge normalmente y satisface
\begin{equation}
 \theta_s(-1/\tau)=(-i\tau)^{12}\theta_{-s}(\tau).
 \label{eq:leech-theta-modular}
\end{equation}
\end{proposition}
\begin{proof}
El valor absoluto de cada sumando es
\(e^{-\pi\operatorname{Im}\tau\|D_s\lambda\|_B^2}\).
Sobre compactos del semiplano superior, la parte imaginaria tiene
una cota inferior positiva, de modo que
\eqref{eq:leech-quadratic-bounds} prueba convergencia normal y
holomorfía. Para \(\tau=it\), con \(t>0\),
\eqref{eq:leech-theta-reciprocity} aplicada en \(t^{-1}\) da
\(\theta_s(i/t)=t^{12}\theta_{-s}(it)\).
Ambos miembros de \eqref{eq:leech-theta-modular} son holomorfos
en \(\mathbb H\); el teorema de identidad prolonga la igualdad
desde el eje imaginario positivo.
\end{proof}

La inversión modular relaciona los parámetros \(s\) y \(-s\).
La periodicidad adicional bajo \(\tau\mapsto\tau+1\) utiliza la
paridad de las normas. La proposición~\ref{prop:leech-nonintegral}
exhibe parámetros arbitrariamente próximos a cero para los que
la integralidad euclídea se pierde. Por ello la colección conserva
la reciprocidad analítica demostrada, mientras una realización
mediante formas modulares escalares o álgebras de vértices
reticulares incorpora además las hipótesis aritméticas
correspondientes. En la realización excepcional inicial, esas
hipótesis sostienen la continuación hacia las álgebras de vértices
y Monstrous Moonshine \cite{HMT-IV}.

\subsection{Forma partida y autodualidad en dimensión cuarenta y ocho}
\label{subsec:leech-split}

La pérdida de integralidad de una proyección euclídea conduce a
considerar ambas proyecciones de manera simultánea. En el espacio
\(E\oplus E\) se introduce la forma bilineal
\begin{equation}
 \mathcal B((x,p),(x',p'))=
       \langle x,p'\rangle_B+\langle p,x'\rangle_B.
 \label{eq:leech-split-form}
\end{equation}
Sus productos emparejan las dos coordenadas. Definimos
\begin{equation}
\begin{split}
 \mathcal D_s&=D_s\oplus D_{-s},\\
 \Gamma_s&=\mathcal D_s(\Lambda\oplus\Lambda)\\
 &=\{(D_s\lambda,D_{-s}\mu):\lambda,\mu\in\Lambda\}.
\end{split}
\label{eq:leech-split-lattice}
\end{equation}

\begin{theorem}[Autodualidad integral para la forma partida]
\label{thm:leech-split}
La forma \(\mathcal B\) tiene signatura \((24,24)\).
Para todo \(s\in\mathbb R\), \(\mathcal D_s\) es una isometría de
\(\mathcal B\), y \(\Gamma_s\) es una red integral, par y autodual
respecto de esta forma.
\end{theorem}
\begin{proof}
Introduciendo
\[
 p_L=\frac{x+p}{\sqrt2},\qquad
 p_R=\frac{x-p}{\sqrt2},
\]
se obtiene
\[
 \mathcal B((x,p),(x,p))=\|p_L\|_B^2-\|p_R\|_B^2.
\]
Como \(B\) es positiva en dimensión \(24\), la signatura es
\((24,24)\). La autoadjunción y reciprocidad de los bloques dan
\[
 \langle D_sx,D_{-s}p'\rangle_B=\langle x,p'\rangle_B,
\]
y análogamente para el segundo producto cruzado. Por ello
\(\mathcal B(\mathcal D_s\xi,\mathcal D_s\eta)=\mathcal B(\xi,\eta)\).

En \(\Gamma_0=\Lambda\oplus\Lambda\), todos los productos cruzados
son enteros y
\[
 \mathcal B((\lambda,\mu),(\lambda,\mu))
   =2\langle\lambda,\mu\rangle_B\in2\mathbb Z.
\]
Esto prueba integralidad y paridad. Si \((x,p)\) pertenece al dual
de \(\Gamma_0\) para \(\mathcal B\), el emparejamiento con
\((\lambda,0)\) y \((0,\mu)\) exige respectivamente
\(p,x\in\Lambda^*=\Lambda\). El dual es, por tanto, \(\Gamma_0\).
La isometría \(\mathcal D_s\) transporta estas tres propiedades
a \(\Gamma_s\).
\end{proof}

La métrica especifica el contenido de la conclusión. La colección
\(\Gamma_s\) conserva la integralidad para los productos cruzados
de \(\mathcal B\), incluso cuando \(\Lambda_s\) pierde integralidad
para \(B\). Como retículos abstractos con forma partida, los
\(\Gamma_s\) son isométricos. Sus dos proyecciones euclídeas,
consideradas con las métricas positivas respectivas, varían con
el parámetro.

\begin{proposition}[Intercambio de las proyecciones recíprocas]
\label{prop:leech-split-involution}
La involución \(\mathcal J(x,p)=(p,x)\) satisface
\[
 \mathcal J\Gamma_s=\Gamma_{-s},\qquad
 \mathcal J\mathcal D_s\mathcal J=\mathcal D_{-s}.
\]
En las coordenadas \((p_L,p_R)\), fija \(p_L\) y cambia el signo
de \(p_R\).
\end{proposition}
\begin{proof}
El intercambio convierte
\((D_s\lambda,D_{-s}\mu)\) en
\((D_{-s}\mu,D_s\lambda)\), que recorre \(\Gamma_{-s}\).
La conjugación de los dos bloques da la segunda identidad.
Finalmente, \(x+p\) se conserva y \(x-p\) cambia de signo.
\end{proof}

La involución expresa una dualidad de la construcción reticular
completa. Su eventual representación física conserva, además de
la forma partida, los operadores y los datos de identificación
propios de esa realización.

\subsection{Energía reticular y equivalencia unitaria}
\label{subsec:leech-energy}

La energía de radio recíproco
\(m^2/R^2+w^2R^2\) intercambia sus dos términos bajo
\((m,w,R)\mapsto(w,m,R^{-1})\), con una conjugación unitaria
que preserva los dominios operatorios \cite{HMT-IV}.
La realización sobre los doce planos marcados emplea
\begin{equation}
 \mathcal E_s(\lambda,\mu)=
       \|D_s\lambda\|_B^2+\|D_{-s}\mu\|_B^2,
 \qquad \lambda,\mu\in\Lambda.
 \label{eq:leech-energy}
\end{equation}
La energía es positiva y cumple
\(\mathcal E_s(\lambda,\mu)=\mathcal E_{-s}(\mu,\lambda)\).

\begin{theorem}[Operador positivo y conjugación de dominios]
\label{thm:leech-energy}
En el espacio de Hilbert \(\ell^2(\Lambda\oplus\Lambda)\), sea
\[
 (H_s\psi)(\lambda,\mu)=
       \mathcal E_s(\lambda,\mu)\psi(\lambda,\mu)
\]
con dominio
\begin{equation}
 \mathcal D(H_s)=
 \left\{\psi\in\ell^2(\Lambda\oplus\Lambda):
 \sum_{\lambda,\mu}\mathcal E_s(\lambda,\mu)^2
             |\psi(\lambda,\mu)|^2<\infty\right\}.
 \label{eq:leech-energy-domain}
\end{equation}
Entonces \(H_s\) es positivo, autoadjunto y de resolvente compacto.
La aplicación \(U\psi(\lambda,\mu)=\psi(\mu,\lambda)\) es unitaria y
\begin{equation}
 U\mathcal D(H_s)=\mathcal D(H_{-s}),\qquad
 UH_sU^{-1}=H_{-s}.
 \label{eq:leech-unitary}
\end{equation}
\end{theorem}
\begin{proof}
La multiplicación por una función real en un conjunto discreto,
con el dominio máximo \eqref{eq:leech-energy-domain}, es
autoadjunta. Puede comprobarse directamente mediante la base
ortonormal de los vectores de soporte unitario: la pertenencia
al dominio del adjunto exige que la sucesión obtenida al
multiplicar por \(\mathcal E_s\) sea cuadrado sumable, exactamente
la misma condición. Además,
\[
 \langle\psi,H_s\psi\rangle
   =\sum_{\lambda,\mu}\mathcal E_s(\lambda,\mu)
                         |\psi(\lambda,\mu)|^2\geq0.
\]
Por \eqref{eq:leech-quadratic-bounds},
\[
 \mathcal E_s(\lambda,\mu)\geq
 e^{-2|s|M}\bigl(\|\lambda\|_B^2+\|\mu\|_B^2\bigr).
\]
Una red tiene finitos puntos en cualquier conjunto acotado; por
ello cada subnivel de \(\mathcal E_s\) es finito. Los coeficientes
diagonales \((1+\mathcal E_s)^{-1}\) tienden a cero fuera de
conjuntos finitos, de modo que \((H_s+I)^{-1}\) es límite en norma
de operadores de rango finito y es compacto.

El intercambio de los dos índices conserva la norma de
\(\ell^2(\Lambda\oplus\Lambda)\), por lo que \(U\) es unitario
y \(U^{-1}=U\). La igualdad de energías al invertir \(s\) transforma
la condición de dominio de \(H_s\) en la de \(H_{-s}\), y la
evaluación sobre cada par \((\lambda,\mu)\) prueba
\eqref{eq:leech-unitary}.
\end{proof}

Quedan así articuladas tres conservaciones de distinto dominio.
En los planos euclídeos, el flujo mantiene covolumen y simetría
ternaria y relaciona cada red con el dual del parámetro opuesto.
En el espacio duplicado, los productos cruzados conservan una
red integral, par y autodual. En la realización operatoria, el
intercambio recíproco conserva la energía mediante una equivalencia
unitaria de operadores con sus dominios. Las fórmulas theta
expresan la misma reciprocidad en la suma convergente de las
longitudes de toda la red.
